\section{The same prior mechanism on an arbitrary fixed finite sphere subset}
\label{sec:arbitrary-finite-domain}
The explicit geometry of $\mathcal U_D$ is needed for a quantitative lower bound on the number of inputs. The maximum-amplitude mechanism itself is more general.

\begin{lemma}[Strict positivity of the native finite-set kernel]
\label{lem:finite-kernel-positive}
Let $\mathcal X$ be any nonempty finite set of distinct points on $\mathbb S^{d-1}$. The native first-layer kernel
\[
 K(x,x')=\E_A\tanh(A\cdot(1,x))\tanh(A\cdot(1,x')),
 \qquad A\sim\mathcal N(0,I/d),
\]
is strictly positive definite on $L^2(\mathcal X)$ for the uniform probability measure.
\end{lemma}
\begin{proof}
The Hermite calculation in Lemma~\ref{lem:sphere-direct-kernel}, with $\nu=0$, shows
\[
 K(x,x')=\sum_{j\text{ odd}}b_j^2
            \left(\frac{1+x\cdot x'}2\right)^j,
 \qquad b_j\ne0\quad\text{for every odd }j.
\]
Let $m=|\mathcal X|$, and choose an odd $j\ge m-1$. The tensor feature map for $(1+x\cdot x')^j$ contains, with positive coefficients, every coordinate monomial of total degree at most $j$. Choose a linear functional $\ell$ taking distinct values on the finitely many points of $\mathcal X$; the exceptional choices form a finite union of proper hyperplanes. Univariate Lagrange interpolation in $\ell(x)$ gives a polynomial of degree at most $m-1$ with any prescribed values on $\mathcal X$. Thus the degree-$j$ tensor features span every finite-set function, and their Gram matrix is positive definite. Its positive coefficient $b_j^2$ proves the claim.
\end{proof}

\begin{corollary}[An arbitrary finite domain with a quantitative kernel parameter]
\label{cor:arbitrary-finite-domain}
Let $\mathcal X\subset\mathbb S^{d-1}$ be finite, let $\Lambda>0$ be the minimum eigenvalue of its native kernel operator, and put
\[
 \Gamma=d+\log(2/\Lambda),\qquad G=\lceil\Gamma\rceil.
\]
There are absolute $A_0,C,H_0>0$ such that, for $H\ge H_0$, $G\le\sqrt H/2$, and every $f:\mathcal X\to[-A,A]$ with $0<A\le A_0$, every fixed $\varepsilon>0$ admits a comparison at all $N\ge C(1+\varepsilon^{-1})^C H^C e^{C\Gamma}$ satisfying
\[
 \E_\rho\max_{x\in\mathcal X}|z_w(x)-f(x)|^2\le\varepsilon^2,
\]
\[
 \frac2N\KL(\rho\|\mu)
 \le4\|f\|_\infty+CA^2+C\frac{\Gamma^2}{H}
       +\frac{2\log2}{N}.
\]
The finite-set entropy lower bound in Corollary~\ref{cor:sup-entropy} applies to the same domain. In particular, every fixed finite sphere subset eventually admits the same leading maximum-amplitude cost as depth and width grow.
\end{corollary}
\begin{proof}
Every first-layer preactivation has variance $2/d$, so its squared-output mean is $b_1=F(2/d)\ge c/d$ for an absolute $c>0$; the bound is uniform also for $d=1$. The kernel operator has trace $b_1\le1$, hence $\Lambda\le1/m$ and $\log m\le\log(1/\Lambda)\le\Gamma$. Lemma~\ref{lem:unit-seed} supplies a seed of size $r_1=\sqrt\Lambda/2$, with $r_1\ge e^{-C\Gamma}$ and $r_1\le\sqrt{b_1}/2$.

Apply Lemma~\ref{lem:spike-prefix} with the complexity parameter $G$ in place of $D$. Its hypotheses hold since $b_1\ge c/(G+1)$ and $\log(1/r_1)=O(G)$. The reverse-tail, padding, and intermediate-value arguments are identical to the proof of Theorem~\ref{thm:canonical-function-memorizer}. Finally $\log m+\log(1/\Lambda)\le2\Gamma=O(\sqrt H)$, so the weighted stability proof of Theorem~\ref{thm:canonical-polynomial-width}, with $G$ in place of $D$, gives the displayed polynomial-in-$H$ threshold with factor $e^{C\Gamma}$. Replacing $G$ by $\Gamma$ changes only absolute constants.
\end{proof}
